\DocumentMetadata{
  pdfversion=2.0,
  pdfstandard=ua-2,
  lang=en-US,
  tagging=on,
  tagging-setup={math/setup=mathml-SE,tabsorder=structure}
}
\documentclass[11pt,letterpaper]{article}
\usepackage[margin=0.68in,includefoot]{geometry}
\usepackage{newcomputermodern}
\usepackage{amsmath,amscd,mathtools}
\usepackage{tikz,adjustbox}
\providecommand{\square}{\mdlgwhtsquare}
\usepackage[hidelinks]{hyperref}
\hypersetup{
  pdftitle={Topology First-Year Exam, January 2023, Part 1},
  pdfauthor={University of Florida Department of Mathematics},
  pdfsubject={Graduate mathematics examination},
  pdfdisplaydoctitle=true,
  bookmarksopen=true
}
\setcounter{secnumdepth}{0}
\setlength{\parindent}{0pt}
\setlength{\parskip}{0.28em}
\setlength{\emergencystretch}{3em}
\setlength{\leftmargini}{2.15em}
\setlength{\leftmarginii}{2.65em}
\setlength{\leftmarginiii}{3em}
\renewcommand{\labelenumi}{\textbf{\theenumi.}}
\pagestyle{empty}
\raggedbottom
\allowdisplaybreaks
\newenvironment{examproblems}[1]{%
  \begin{enumerate}%
  \setcounter{enumi}{#1}%
  \setlength{\topsep}{0.35em}%
  \setlength{\partopsep}{0pt}%
  \setlength{\itemsep}{0.48em}%
  \setlength{\parsep}{0.18em}%
}{\end{enumerate}}
\newenvironment{examsubparts}{%
  \begin{enumerate}%
  \setlength{\topsep}{0.18em}%
  \setlength{\partopsep}{0pt}%
  \setlength{\itemsep}{0.16em}%
  \setlength{\parsep}{0.1em}%
}{\end{enumerate}}
\newenvironment{examinstructions}{%
  \par\begingroup\itshape
}{%
  \par\endgroup\vspace{0.18em}
}
\makeatletter
\renewcommand\section{\@startsection{section}{1}{0pt}%
  {-1.25ex plus -.25ex minus -.1ex}%
  {0.55ex plus .1ex}%
  {\normalfont\large\bfseries}}
\renewcommand{\@maketitle}{%
  \newpage\null\vskip -1em%
  {\footnotesize\bfseries
    \href{https://www.ufl.edu/}{University of Florida}, \href{https://math.ufl.edu/}{Department of Mathematics}\par}%
  \vskip -0.5em%
  \tagstructbegin{tag=Title}%
    {\LARGE\bfseries\href{https://gma.math.ufl.edu/exams/topology-first-year/}{Topology First-Year Exam}, January 2023, Part 1\par}%
  \tagstructend%
  \vskip 0.25em%
  \tagmcbegin{artifact}%
    \hrule height 1pt%
  \tagmcend%
  \vskip 1em%
}
\makeatother
\title{Topology First-Year Exam, January 2023, Part 1}
\author{}
\date{}
\begin{document}
\maketitle
\thispagestyle{empty}
\begin{examinstructions}
Work the following problems and show all work. Support all statements to the best of your ability.
\end{examinstructions}
\begin{examproblems}{0}
\item
Show that every map \(f:S^2\to T\) of the 2-sphere to the torus~\(T\) is null-homotopic.
\item
Show that the torus~\(T\) is not homeomorphic to the surface of genus 2, \(M_2=T\mathbin{\#}T\).
\item
Show that for every continuous map \(f:S^1\to S^1\) there is \(n\in\mathbb{Z}\) such that~\(f\) is homotopic to \(z^n:S^1\to S^1\).
\item
Let \(X=[-1,1]\times\{0\}\cup\{0\}\times[-1,1]\subset\mathbb{R}\times\mathbb{R}\). Show that every continuous map \(f:X\to X\) has a fixed point.
\item
Let~\(X\), \(Y\), and~\(Z\) denote the~\(x\)-axis, the~\(y\)-axis, and the~\(z\)-axis in \(\mathbb{R}^3\). Is \(\pi_1(A)\) abelian where
\begin{examsubparts}
\item[\textnormal{(a)}]
\(A=\mathbb{R}^3\setminus X\)?
\item[\textnormal{(b)}]
\(A=\mathbb{R}^3\setminus(X\cup Y)\)?
\item[\textnormal{(c)}]
\(A=\mathbb{R}^3\setminus(X\cup Y\cup Z)\)?
\end{examsubparts}
\end{examproblems}
\begin{examinstructions}
Answer the following with complete definitions or statements or proofs.
\end{examinstructions}
\begin{examproblems}{5}
\item
Show that \(S^2\) is not homeomorphic to \(S^3\).
\item
State the Borsuk–Ulam Theorem for \(S^2\).
\item
Let \(I=[0,1]\) and let \(X=I^I\) be given the product topology. Is the space~\(X\)
\begin{examsubparts}
\item[\textnormal{(a)}]
compact?
\item[\textnormal{(b)}]
metrizable?
\item[\textnormal{(c)}]
separable?
\end{examsubparts}
\item
State the Seifert–van Kampen Theorem. Can it be used to compute \(\pi_1(S^1)\)?
\item
This problem has two parts.
\begin{examsubparts}
\item[\textnormal{(a)}]
What are completely regular spaces?
\item[\textnormal{(b)}]
Define the Stone–Čech compactification \(\beta X\) of~\(X\).
\end{examsubparts}
\end{examproblems}
\end{document}
